% ============================================================================
% Chapter 5 --- Architectural principles
% ----------------------------------------------------------------------------
% Purpose : establish the principles common to all reference-architecture
%           chapters: interfaces over products, replaceability as a
%           first-class property, vendor-neutral by construction.
%           Freezes the standards landscape table (single source of truth
%           for Ch. 6--14).
% Page target: ~4 pages.
% Dependencies: Ch. 3 (sovereignty grid), Ch. 4 (self-assessment).
% ----------------------------------------------------------------------------
% Decisions registered in _coherence-EN.md:
%   D5.1 nine functional layers grouped in three planes (Subject/Execution/Control)
%   D5.2 standards table is normative; products in Ch. 6--14 are illustrative
%   D5.3 orchestration plane has no single standard (acknowledged honestly)
%   D5.4 replaceability defined operationally via Ch. 16 conformance tests
% ============================================================================

\chapter{Architectural principles}
\label{ch:architectural-principles}

\begin{caveatbox}[Dependencies]
\noindent This chapter rests on the sovereignty grid of
\trchap{ch:sovereignty-grid}{} --- in particular on the technical
dimension. It establishes the principles that the layer chapters
(Chapters~6--14 of the Technical Reference) will instantiate, and
freezes the standards landscape that those chapters reference.
\end{caveatbox}

The reference architecture developed in Part~II is governed by a small
number of principles. They are stated here so that the per-layer
chapters need only apply them, and so that any deviation in those
chapters is visible against an established baseline.

%-----------------------------------------------------------------------------
\section{Interfaces over products}
\label{sec:ap-interfaces}

The normative content of this architecture is the set of interface
contracts that connect the layers. The reference instantiations
proposed in the per-layer chapters --- open-source implementations of
identity providers, policy engines, network access servers, cloud
infrastructure managers, and telemetry pipelines, among others --- are
illustrative: they demonstrate that the contracts are realisable on
open foundations, not that the institution must adopt these specific
products.

Three consequences follow from this primacy:

\begin{itemize}[leftmargin=1.5em,nosep]
  \item \textbf{Procurement language} references standards, not
    products. An \gls{rfp} that asks for ``identity provider compatible
    with OIDC Core 1.0, SCIM 2.0, and federation metadata in
    accordance with the applicable research-and-education identity
    federation technical profile'' admits multiple
    compliant suppliers; an \gls{rfp} that asks for a named
    proprietary identity suite with its vendor-specific conditional
    access feature admits one.
  \item \textbf{Component evaluation} consists in verifying that the
    standard is spoken correctly, not that the supplier promises a
    feature. The conformance test suite of
    \trchap{ch:conformance}{Conformance and replaceability tests}
    operationalises this verification.
  \item \textbf{Architectural decisions} are recorded against the
    contract, not against the product. A decision to use OIDC for
    \gls{sso} survives the replacement of the identity provider; a
    decision to use a specific vendor's proprietary conditional
    access expression \texttt{X} does not.
\end{itemize}

This principle is the operational expression of the technical
sovereignty dimension: an institution that controls the contracts
controls the architecture, regardless of the suppliers that currently
realise those contracts.

%-----------------------------------------------------------------------------
\section{Functional necessity is independent of scale}
\label{sec:ap-necessity}

The interface contracts of \S\ref{sec:ap-interfaces} are not only the
normative content of the architecture; they enumerate the
\emph{minimal set of functions that any institution offering a digital
infrastructure must discharge}, irrespective of its size, budget, or
research intensity. The moment an institution authenticates users,
holds personal or sensitive data, attaches devices to a network, and
operates services, it has already incurred these functions. The only
question is whether it discharges them coherently and sovereignly, or
implicitly and dependently.

\paragraph{Several of these functions are legally compelled, without a
size threshold.} The \gls{gdpr}~\cite{gdpr} applies to any institution
processing personal data, with no exemption by size; the
data-protection, segmentation and audit functions are therefore
obligations rather than aspirations at every scale. Where the
institution is classified as an essential or important entity, the
\gls{nis2}~\cite{nis2} obligations apply in addition. The minimal set
is thus anchored in binding instruments that do not vary with the size
of the institution.

\paragraph{What is size-invariant, and what scales.} What does
\emph{not} vary with scale is the set of functional contracts and the
legally compelled protections above. What \emph{does} scale is the
\emph{instantiation} (the specific products that realise a contract),
the \emph{sizing} (\trchap{ch:sizing}{}, given as orders of magnitude
for \loc{institution-type}, with a down-scaled profile in
\S\ref{sec:sz-downscale}), and the \emph{optional components} justified
only by particular needs --- a dedicated Science~DMZ, multi-site
federation --- which a smaller or teaching-led institution may
legitimately defer or mutualise. The architecture is the invariant;
the realisation is the variable.

\paragraph{Consequence for the ``we are not a flagship'' objection.}
The objection that an institution ``cannot afford what the leading
universities do'' misidentifies the commitment. The architecture asks
no institution to match a flagship's budget or estate; it observes
that the functions are already owed the moment digital services are
offered, and provides the minimal coherent and sovereign way to
discharge them. The alternative to adopting the architecture is not
\emph{not discharging these functions} --- it is discharging the same
functions incoherently and dependently.

%-----------------------------------------------------------------------------
\section{Replaceability as a first-class property}
\label{sec:ap-replaceability}

A component is \emph{replaceable} when it can be substituted by
another one speaking the same interface contract, within an agreed
time window, with the institution's data and configuration preserved.
Replaceability is the property that makes the architecture survive
supplier change, regulatory change, and the obsolescence of any
particular product.

Four conditions, jointly, define replaceability.

\begin{description}[leftmargin=1.5em,style=nextline,nosep]
  \item[Standard interfaces.] The component speaks the relevant
    interface standard without captive extensions in the critical
    path. Optional extensions are admissible if they can be turned
    off without loss of essential function.
  \item[Open state and configuration formats.] The component's
    operational state and configuration are expressed in formats that
    can be read and produced by parties other than the supplier.
  \item[Documented exit-data export.] A procedure exists to extract
    the institution's data in a documented open format within an
    agreed time window, without requiring the supplier's discretionary
    cooperation.
  \item[No hidden dependencies.] The component does not depend on
    supplier-specific operators, controllers or sidecars whose absence
    would silently degrade behaviour.
\end{description}

Replaceability is not the same as \emph{having already been replaced}.
The conformance suite of
\trchap{ch:conformance}{} includes an \emph{exit rehearsal} category:
the institution periodically exercises the replacement procedure on a
representative subset, to confirm that the four conditions still
hold.

%-----------------------------------------------------------------------------
\section{Standards landscape}
\label{sec:ap-standards}

Table~\ref{tab:ap-standards} lists the standards on which the
reference architecture rests. The list is ordered by layer, then by
the criticality of the standard within that layer. It is the canonical
reference for Chapters~6--14: any layer chapter that introduces a
standard not listed here either introduces a documented exception or
extends this table.

\begin{table}[H]
\centering\small
\caption{Standards landscape of the reference architecture. Status
indicates the maturity of the standard relative to the architecture's
expectations: \emph{Stable} (long-established, broadly implemented),
\emph{Mature} (established but evolving), \emph{Emerging} (recent,
adopted by the architecture as forward-looking).}
\label{tab:ap-standards}
\begin{tabularx}{\textwidth}{%
  >{\hsize=0.55\hsize\linewidth=\hsize\bfseries\raggedright\arraybackslash}X%
  >{\hsize=0.95\hsize\linewidth=\hsize\raggedright\arraybackslash}X%
  >{\hsize=1.05\hsize\linewidth=\hsize\raggedright\arraybackslash}X%
  >{\hsize=0.45\hsize\linewidth=\hsize\raggedright\arraybackslash}X%
}
\toprule
Layer & Standard & Role in the architecture & Status \\
\midrule
\multirow{5}{*}{Identity}
  & OIDC Core~1.0~\cite{oidc-core-1.0}      & Federated authentication, identity assertion & Stable \\
  & SAML~2.0~\cite{saml-core-2.0}           & Federated authentication (legacy compatibility) & Stable \\
  & SCIM~2.0~\cite{rfc-7643-scim-core,rfc-7644-scim-protocol} & Identity provisioning across systems & Stable \\
  & OAuth~2.1~\cite{oauth-2-1}              & Authorisation framework underpinning OIDC & Emerging \\
  & OpenID Federation  & Trust establishment between federations & Emerging \\
\addlinespace
\multirow{2}{*}{Policy}
  & OPA REST API~\cite{opa-project}         & Policy decision point query interface (de facto project API) & Mature \\
  & Cedar evaluation API~\cite{cedar-project} & Alternative policy decision interface (project specification) & Emerging \\
\addlinespace
\multirow{5}{*}{Network}
  & RADIUS~\cite{rfc-2865-radius,rfc-2866-radius-acct} & Network access authentication and accounting & Stable \\
  & NETCONF~\cite{rfc-6241-netconf}\,/\,YANG~\cite{rfc-7950-yang} & Network device configuration and modelling & Stable \\
  & BGP / EVPN         & Routing and overlay signalling & Stable \\
  & CNI~\cite{cni-spec} / eBPF        & Kubernetes network plugin model (CNCF project specification) & Mature \\
  & IPFIX~\cite{rfc-7011-ipfix}             & Flow record export & Stable \\
\addlinespace
\multirow{4}{*}{Compute}
  & Kubernetes API~\cite{kubernetes}        & Container orchestration & Stable \\
  & OpenStack API~\cite{openstack}          & Infrastructure as a Service & Stable \\
  & OCI image format~\cite{oci-image-spec}  & Container image packaging & Stable \\
  & SLSA~\cite{slsa-spec}                   & Software supply-chain provenance & Mature \\
\addlinespace
\multirow{3}{*}{Storage}
  & S3 API             & Object storage interface (de facto standard) & Stable \\
  & KMIP~\cite{kmip-oasis}                  & Key management interoperability & Stable \\
  & NFS v4.1~\cite{rfc-8881-nfs41}          & Network file system & Stable \\
\addlinespace
\multirow{3}{*}{Endpoint}
  & DDM / DEP          & Apple declarative device management (vendor specification) & Stable \\
  & MS-MDM             & Microsoft device management (published vendor specification) & Stable \\
  & Declarative Linux config~\cite{osquery,ansible} & Linux endpoint posture and configuration (project tooling, no standards-body specification) & Mature \\
\addlinespace
\multirow{3}{*}{Observability}
  & OpenTelemetry~\cite{opentelemetry-spec} & Telemetry collection (traces, metrics, logs) & Stable \\
  & CloudEvents~\cite{cloudevents-spec}     & Event format for cross-system signalling & Mature \\
  & syslog~\cite{rfc-5424-syslog}           & Legacy log transport & Stable \\
\addlinespace
\multirow{4}{*}{Crypto}
  & PKCS\#11~\cite{pkcs11-oasis}            & Hardware security module interface & Stable \\
  & ACME~\cite{rfc-8555-acme}               & Automated certificate issuance & Stable \\
  & KMIP~\cite{kmip-oasis}                  & Key management interoperability (cross-listed) & Stable \\
  & OpenID Federation  & Trust hierarchy for federated entities & Emerging \\
\addlinespace
Orchestration & (no single standard --- see \S\ref{sec:ap-standards} note below) & Cross-layer command issuance after a TAC activation & --- \\
\bottomrule
\end{tabularx}
\end{table}

\paragraph{The orchestration gap.} The orchestration plane --- the
``glue'' that turns a TAC activation decision into the cross-layer
commands that realise it --- has no single standard at the time of
writing. Several patterns are converging (event-driven workflow
engines, declarative reconcilers, GitOps repositories), but no
specification consolidates them into a portable interface. This is
the architecture's most honest gap; \trchap{ch:orchestration}{The
orchestration plane} discusses how to bound the consequences and how
the institution can avoid creating a new captive dependency in this
layer.

\paragraph{The policy-interface choice.} Unlike the orchestration
plane, the policy plane does have a standards-body specification ---
OASIS XACML --- which the architecture deliberately does not adopt:
the de facto project APIs it retains (OPA REST, Cedar) are the
interfaces its reference instantiations implement, and the table
records their project status accordingly.

\paragraph{Standards versus reference profiles.} Several standards
listed above (OIDC, SAML2, OPA, OpenTelemetry) admit broad
configuration latitude. Where the architecture relies on a specific
profile of a standard --- for instance, the applicable
research-and-education identity federation technical profile for SAML2
federation, or the OPA bundle distribution pattern --- the
profile is named in the relevant layer chapter and treated as a
co-normative element.

%-----------------------------------------------------------------------------
\section{Layer decomposition}
\label{sec:ap-layers}

The nine functional layers of the architecture are organised in three
planes that describe what each layer governs.

\begin{description}[leftmargin=1.5em,style=nextline,nosep]
  \item[Subject and trust plane.] Identity (\trchap{ch:identity}{}),
    Endpoint (\trchap{ch:endpoint}{}), Cryptographic services
    (\trchap{ch:crypto}{}). These layers together define \emph{who}
    is acting, \emph{from where}, and \emph{with what cryptographic
    assurances}.
  \item[Execution plane.] Network (\trchap{ch:network}{}), Compute
    (\trchap{ch:compute}{}), Storage (\trchap{ch:storage}{}). These
    layers together define \emph{where} the activity runs and on
    \emph{which data}.
  \item[Control and governance plane.] Policy
    (\trchap{ch:policy}{}), Observability
    (\trchap{ch:observability}{}), Orchestration
    (\trchap{ch:orchestration}{}). These layers together define
    \emph{how} decisions are made, recorded, and executed across the
    other planes.
\end{description}

The dependency structure has two regular features. First, the subject
and trust plane is read by the control plane (Identity attributes
parameterise Policy decisions; Endpoint posture conditions Policy and
Network attachment; Crypto underpins all assertions). Second, the
control plane writes to the execution plane (Policy decisions are
enforced at Network, Compute, and Storage; Orchestration realises
them across the three; Observability collects evidence of all
crossings). The reference architecture is the realisation of these two
flows.

\begin{figure}[H]
\centering
\begin{tikzpicture}[
    every node/.style = {font=\footnotesize},
    layer/.style = {rectangle, rounded corners=2pt,
                    draw=NavyBlue!70!black, fill=NavyBlue!5,
                    minimum width=26mm, minimum height=10mm,
                    line width=0.4pt, align=center},
    plane/.style = {rectangle, rounded corners=4pt, draw=Slate!50,
                    line width=0.4pt, dashed, inner sep=5pt},
    arrow/.style = {-{Stealth[length=2mm]}, draw=NavyBlue!50!black,
                    line width=0.4pt},
    flowarrow/.style = {-{Stealth[length=2.5mm]}, draw=Crimson!70!black,
                        line width=0.6pt}
]
% Subject and Trust plane (top)
\node[layer] (id) at ( 0.0, 0.0) {Identity\\ \tiny (Ch.~6)};
\node[layer] (ep) at ( 3.2, 0.0) {Endpoint\\ \tiny (Ch.~11)};
\node[layer] (cr) at ( 6.4, 0.0) {Crypto\\ \tiny (Ch.~13)};
\node[plane, fit=(id)(ep)(cr), label={[font=\scriptsize\bfseries,
      color=Slate!80!black]above:Subject and Trust plane}] {};

% Execution plane (middle)
\node[layer] (nw) at ( 0.0,-2.8) {Network\\ \tiny (Ch.~8)};
\node[layer] (co) at ( 3.2,-2.8) {Compute\\ \tiny (Ch.~9)};
\node[layer] (st) at ( 6.4,-2.8) {Storage\\ \tiny (Ch.~10)};
\node[plane, fit=(nw)(co)(st), label={[font=\scriptsize\bfseries,
      color=Slate!80!black]left:Execution plane}] {};

% Control and Governance plane (bottom)
\node[layer] (po) at ( 0.0,-5.6) {Policy\\ \tiny (Ch.~7)};
\node[layer] (ob) at ( 3.2,-5.6) {Observability\\ \tiny (Ch.~12)};
\node[layer] (or_) at ( 6.4,-5.6) {Orchestration\\ \tiny (Ch.~14)};
\node[plane, fit=(po)(ob)(or_), label={[font=\scriptsize\bfseries,
      color=Slate!80!black]below:Control and Governance plane}] {};

% Flow 1 (Crimson): Subject/Trust -> Control plane
\draw[flowarrow] (id.south) to[bend right=10] (po.north west);
\draw[flowarrow] (ep.south) to[bend right=8]  (po.north);
\draw[flowarrow] (cr.south) to[bend right=10] (or_.north);

% Flow 2 (Crimson): Control plane -> Execution plane
\draw[flowarrow] (po.north east) to[bend left=12] (nw.south);
\draw[flowarrow] (po.north east) to[bend left=8]  (co.south);
\draw[flowarrow] (po.north east) to[bend left=8]  (st.south);
\draw[flowarrow] (or_.north west) to[bend right=15] (co.south east);

% Observability collects (lighter arrows)
\draw[arrow, dashed] (nw.south west) to[bend right=10] (ob.north west);
\draw[arrow, dashed] (co.south) -- (ob.north);
\draw[arrow, dashed] (st.south) to[bend left=10]  (ob.north east);

% Legend
\node[font=\scriptsize, anchor=west] at (8.0,-1.0) {%
  \textcolor{Crimson!70!black}{\textbf{---}} principal flows};
\node[font=\scriptsize, anchor=west] at (8.0,-1.6) {%
  \textcolor{NavyBlue!50!black}{\textbf{- - -}} audit collection};
\end{tikzpicture}
\caption{Layer decomposition of the reference architecture: nine
functional layers organised in three planes. The principal
dependency flows (Crimson) are: Subject and Trust feeds the
Control plane; the Control plane writes to the Execution plane.
Observability (dashed) collects events from every layer.}
\label{fig:ap-layers}
\end{figure}

%-----------------------------------------------------------------------------
\section{Availability and failure semantics of the mediation planes}
\label{sec:ap-availability}

Per-session policy decision makes the identity and policy planes
mediation points for the whole institution: left unspecified, their
unavailability would be everyone's unavailability. The architecture
therefore treats availability and failure semantics as first-class
design obligations of the mediation planes, on a par with
replaceability.

\begin{description}[leftmargin=1.5em,style=nextline,nosep]
  \item[Redundant deployment per critical plane.] The identity
    provider, the policy decision point, and the certificate and
    secrets services run in redundant deployments --- active-active,
    or active-passive with a documented and rehearsed failover ---
    with an availability target declared per plane and measured
    through the observability layer. The reference components
    support clustered operation; the sizing envelopes of
    \trchap{ch:sizing}{} assume it.
  \item[A policy-plane outage must not interrupt established work.]
    Policy enforcement points cache decisions for the lifetime of
    the activation token: an outage of the decision point prevents
    \emph{new} activations but does not invalidate sessions already
    activated. Token lifetime is the institution's tunable
    trade-off between revocation latency and outage tolerance,
    declared per zone class.
  \item[Fail-mode policy per zone class.] Each zone's security
    profile declares its fail mode for new attachments during a
    mediation-plane outage: \emph{fail-closed} for administrative
    and sensitive zones; \emph{fail-degraded with audit} where the
    availability of teaching or BYOD connectivity outweighs the
    interception of new sessions. Absent a declared choice, the
    default is fail-closed.
  \item[The commercial-SLA comparison.] A managed identity or
    policy service answers the same single-point-of-failure risk
    with a contractual remedy after the fact; this architecture
    answers it with engineered redundancy under institutional
    control and with declared degradation semantics that the
    institution can design, rehearse, and audit.
\end{description}

%-----------------------------------------------------------------------------
\section{What this reference architecture is and is not}
\label{sec:ap-scope}

The clarity of subsequent chapters depends on a shared understanding
of what this architecture commits to and what it does not.

\paragraph{It is.} A contract-driven architecture, organised in nine
layers and three planes, whose realisability is demonstrated on open
foundations and which integrates proprietary components that respect
the contracts. Its normative content is the standards landscape of
\S\ref{sec:ap-standards} and the four conditions of replaceability of
\S\ref{sec:ap-replaceability}.

\paragraph{It is not.} A product recommendation: the per-layer
chapters present reference instantiations and alternatives, but the
choice belongs to the institution and is informed by the sovereignty
grid (\trchap{ch:sovereignty-grid}{}), not by this Reference. Nor is
it a deployment specification: sizing
(\trchap{ch:sizing}{}) is given as orders of magnitude for \loc{institution-type}. Nor is it a security certification: the conformance
suite (\trchap{ch:conformance}{}) verifies interface fidelity and
replaceability, not absence of vulnerability.

\paragraph{Why the boundary matters for international
generalisation.} The architecture is designed to transpose across
institutions and jurisdictions because its commitments are at the
level of contracts and properties, not at the level of products and
configurations. An institution in a different regulatory context, or
with a different existing estate, instantiates the same contracts
with different products without leaving the architecture. This is the
condition under which the network of universities anticipated in
phase~2 can share an architectural baseline despite operating in very
different conditions.
